% Merged 2026: folds ai_educators_howaiworks/_stateofplay/_taskvsjob/_programstructure
% into one file (they were 4 small standalone files, 15-72 lines each). Section boundaries
% are marked inline with \section so this single \input still produces 4 separate sections
% in the deck, exactly as before the merge.

\section[HowAI]{How AI Actually Works}

%%%%%%%%%%%%%%%%%%%%%%%%%%%%%%%%%%%%%%%%%%%%%%%%%%%%%%%%%%%%%%%%%%%%%%%%%%%%%%%%%%
\begin{frame}[fragile]\frametitle{}
\begin{center}
{\Large How AI Actually Works}

\medskip
{\small (Just enough to demystify it, not a computer science course)}
\end{center}
\end{frame}

%%%%%%%%%%%%%%%%%%%%%%%%%%%%%%%%%%%%%%%%%%%%%%%%%%%%%%%%%%%
\begin{frame}[fragile]\frametitle{What Is a Language Model?}
\begin{itemize}
\item While typing an SMS, have you noticed it suggests the next word?
\item Those suggestions are not random: the phone has learned what word usually comes next from millions of past sentences.
\item Do this once, word by word, and you get a suggestion. Do it again and again, and you get a full sentence.
\item That is a \textbf{Language Model}, predicting the next word. Chained together, it writes whole paragraphs: a \textbf{Generative Model}.
\end{itemize}
\begin{center}
\includegraphics[width=0.4\linewidth,keepaspectratio]{chatgpt34}
\end{center}
\end{frame}

%%%%%%%%%%%%%%%%%%%%%%%%%%%%%%%%%%%%%%%%%%%%%%%%%%%%%%%%%%%
\begin{frame}[fragile]\frametitle{Why Are They Called ``Large''?}
\begin{itemize}
\item ChatGPT works on exactly this idea, just at an enormous scale: trained on a large fraction of everything ever written on the internet.
\item That is why the phrase is \textbf{Large Language Model (LLM)}: large training data, large number of internal parameters.
\item More data and more parameters mean better predictions, but the underlying idea never changes: predict the next word, one at a time.
\end{itemize}
\begin{center}
\includegraphics[width=0.5\linewidth,keepaspectratio]{chatgpt31}
\end{center}
\end{frame}

%%%%%%%%%%%%%%%%%%%%%%%%%%%%%%%%%%%%%%%%%%%%%%%%%%%%%%%%%%%
\begin{frame}[fragile]\frametitle{The Big Picture: AI, ML, and Deep Learning}
\begin{center}
\includegraphics[width=0.55\linewidth,keepaspectratio]{ai1}
\end{center}
\begin{itemize}
\item \textbf{AI}: any machine doing something that looks intelligent
\item \textbf{Machine Learning (ML)}: a machine that learns the pattern from data, instead of being told the rule directly
\item \textbf{Deep Learning (DL)}: ML using many-layered neural networks, powerful enough to learn something as complex as language
\end{itemize}
{\tiny (Ref: https://blogs.nvidia.com/blog/2016/07/29/whats-difference-artificial-intelligence-machine-learning-deep-learning-ai/)}
\end{frame}

\section[State]{Where AI in Education Stands}

%%%%%%%%%%%%%%%%%%%%%%%%%%%%%%%%%%%%%%%%%%%%%%%%%%%%%%%%%%%%%%%%%%%%%%%%%%%%%%%%%%
\begin{frame}[fragile]\frametitle{}
\begin{center}
{\Large Where AI in Education Actually Stands}
\end{center}
\end{frame}

%%%%%%%%%%%%%%%%%%%%%%%%%%%%%%%%%%%%%%%%%%%%%%%%%%%
\begin{frame}[fragile]\frametitle{What Teachers Are Actually Using It For}
\begin{itemize}
\item Teacher AI adoption has grown fast worldwide over the last two years, surveys everywhere point the same direction.
\item The most common uses are unglamorous but time-consuming: gathering research and content, drafting lesson plans, summarizing information, generating classroom materials.
\item \textbf{Notice what all four have in common:} none of them is ``teach the class'' or ``know the student.''
\end{itemize}
\end{frame}

%%%%%%%%%%%%%%%%%%%%%%%%%%%%%%%%%%%%%%%%%%%%%%%%%%%
\begin{frame}[fragile]\frametitle{Opportunities AI Opens Up}
\begin{itemize}
\item \textbf{Differentiated learning, without extra prep:} the same worksheet re-explained at three different difficulty levels, generated in the time it used to take to write one.
\item \textbf{Patient, always-on practice for struggling students:} a student who needs the same concept explained five different ways gets it without feeling like a burden on the teacher or the class.
\item \textbf{A bridge for the language gap:} explaining or translating a difficult English textbook concept into Marathi (or any regional language) on the spot, so understanding does not wait for a bilingual teacher to be free.
\item \textbf{Reclaimed time:} the hours saved on drafting and admin (previous slide) are hours that can go back into actually noticing and helping individual students.
\end{itemize}
\end{frame}

\section[TaskvJob]{Will AI Replace Teachers?}

%%%%%%%%%%%%%%%%%%%%%%%%%%%%%%%%%%%%%%%%%%%%%%%%%%%%%%%%%%%%%%%%%%%%%%%%%%%%%%%%%%
\begin{frame}[fragile]\frametitle{}
\begin{center}
{\Large Will AI Replace Teachers?}
\end{center}
\end{frame}

%%%%%%%%%%%%%%%%%%%%%%%%%%%%%%%%%%%%%%%%%%%%%%%%%%%
\begin{frame}[fragile]\frametitle{2016: A Very Confident Prediction}
\begin{itemize}
\item Every major AI advance since the 1980s has come with the same headline: ``this job is over.'' Radiology got there first, and loudest, a full decade before ChatGPT existed, so what actually happened to it is the single most useful data point we have for thinking about teaching.
\item Geoffrey Hinton, the ``Godfather of AI'' and a Nobel laureate, told a room of doctors in 2016:
\end{itemize}
\begin{alertblock}{Hinton, 2016}
``We should stop training radiologists now. It's just completely obvious that within five years, deep learning is going to do better than radiologists.''
\end{alertblock}
\begin{itemize}
\item He later said, to the New York Times, that he had spoken too broadly.
\end{itemize}
{\tiny (Ref: New York Times, revisited in \emph{radiologybusiness.com}, 2026.)}
\end{frame}

%%%%%%%%%%%%%%%%%%%%%%%%%%%%%%%%%%%%%%%%%%%%%%%%%%%
\begin{frame}[fragile]\frametitle{What Actually Happened}
\begin{itemize}
\item Radiologists did not disappear. Demand for them grew, and it is still growing.
\item AI got very good at one task inside the job: reading a scan.
\item Everything else, consulting with physicians, explaining findings to a worried patient, deciding what to biopsy, stayed entirely human.
\item \textbf{The lesson:} a job survives as long as most of its tasks resist automation, even when one high-profile task falls.
\end{itemize}
\end{frame}

%%%%%%%%%%%%%%%%%%%%%%%%%%%%%%%%%%%%%%%%%%%%%%%%%%%
\begin{frame}[fragile]\frametitle{The Same Logic Applies to Teaching}
\begin{itemize}
\item ``Teacher'' is also a bundle of tasks, not one thing.
\item \textbf{Highly automatable:} drafting a first-pass lesson outline, generating quiz questions, summarizing 50 feedback forms, writing a first draft of report comments.
\item \textbf{Not automatable any time soon:} noticing a student is struggling before they say so, managing a room of 30 restless 12-year-olds, a difficult conversation with a parent, being the adult a kid trusts.
\item AI is quietly eating the first list. Nobody is close to touching the second.
\end{itemize}
\end{frame}

%%%%%%%%%%%%%%%%%%%%%%%%%%%%%%%%%%%%%%%%%%%%%%%%%%%
\begin{frame}[fragile]\frametitle{Quick Check: Task or Job?}
\begin{itemize}
\item Nvidia's CEO Jensen Huang put it this way in 2026: \emph{``AI kills tasks, not jobs.''}
\item For each of these, is AI automating a \textbf{task} within teaching, or threatening the \textbf{job} of teaching itself?
\begin{itemize}
\item Grading 40 multiple-choice quizzes
\item Resolving a conflict between two students at recess
\item Drafting a first pass at a rubric
\item Deciding whether a struggling student needs extra support
\end{itemize}
\end{itemize}
\end{frame}

%%%%%%%%%%%%%%%%%%%%%%%%%%%%%%%%%%%%%%%%%%%%%%%%%%%
\begin{frame}[fragile]\frametitle{Quick Check: Answer}
\begin{itemize}
\item Grading MCQs, drafting a rubric: \textbf{tasks}. Mechanical, pattern-based, already being automated, and that is fine.
\item Resolving a peer conflict, judging who needs support: \textbf{the job}. These need judgment, context, and a relationship built over months, exactly what nobody has figured out how to automate.
\item \textbf{The reframe:} the question is never ``will AI replace teachers,'' it is ``which specific tasks on my plate can I hand off this term.''
\end{itemize}
\end{frame}

\section[Programs]{How These Programs Are Structured}

%%%%%%%%%%%%%%%%%%%%%%%%%%%%%%%%%%%%%%%%%%%%%%%%%%%%%%%%%%%%%%%%%%%%%%%%%%%%%%%%%%
\begin{frame}[fragile]\frametitle{}
\begin{center}
{\Large How ``AI for Teachers'' Programs Are Actually Structured}
\end{center}
\end{frame}

%%%%%%%%%%%%%%%%%%%%%%%%%%%%%%%%%%%%%%%%%%%%%%%%%%%
\begin{frame}[fragile]\frametitle{The Four Pieces, Every Time}
\begin{itemize}
\item Whatever the branding, every serious education-AI program is built from the same four pieces.
\item \textbf{A skills library:} ready-made prompts and workflows for recurring teacher tasks, lesson planning, rubric drafting, differentiated explanations, so nobody starts from a blank prompt box.
\item \textbf{Standards grounding:} lesson plans are checked against your actual curriculum, not the model's generic training knowledge.
\item \textbf{Task automation:} tools that run a recurring job on their own, e.g. a daily review of submitted exit tickets, without you re-prompting every time.
\item \textbf{Privacy guardrails:} clear, explicit rules about what happens to student data.
\end{itemize}
\end{frame}

%%%%%%%%%%%%%%%%%%%%%%%%%%%%%%%%%%%%%%%%%%%%%%%%%%%
\begin{frame}[fragile]\frametitle{What to Look For in Any Tool Your School Considers}
\begin{itemize}
\item Does it know your curriculum, or is it just guessing from general training knowledge?
\item Can it run on autopilot for the boring, repetitive parts, or does every use start from scratch?
\item Does it say clearly, in plain language, what happens to student data?
\item If the answer to any of these is unclear, ask before adopting it school-wide.
\end{itemize}
\end{frame}
